\section*{Conclusion and Consolidated Recommendations}
\addcontentsline{toc}{section}{Conclusion and Consolidated Recommendations}
This report has systematically analyzed and resolved 19 critical issues spanning the entire workflow for training a Curriculum-Transfer Learning Physics-Informed Neural Network (CTL-PINN) for the three-body problem. By establishing definitive, physically-grounded, and computationally robust solutions, this document transforms a set of inconsistent and incomplete instructions into a single, coherent, and actionable implementation blueprint.

The key resolutions have unified the project's approach to foundational physics, data normalization, network architecture, loss function design, and training strategy. The proposed framework now rests on a solid foundation, with clear specifications for computing curriculum metrics like Lyapunov time, a consistent $G=1$ dimensionless unit system, a curriculum-aware schedule for the SIREN $\omega_0$ parameter, a rigorous state management protocol for the ALM, and an efficient, unified data sampling strategy that gracefully handles mixed-resolution data.

For successful implementation, adherence to the consolidated specifications is paramount. The following table provides a final, master summary of all key dynamic hyperparameters, serving as a single source of truth for configuring the training process across all curriculum stages. By following this blueprint, the project is well-positioned to develop a state-of-the-art model capable of learning the complex, chaotic dynamics of the three-body problem with high fidelity and long-term physical plausibility.

\begin{table}[H]
\centering
\caption{Consolidated Hyperparameter Schedule for CTL-PINN Training}
\label{tab:hyperparameter_schedule}
\renewcommand{\arraystretch}{1.2}
\begin{tabular}{p{0.18\linewidth} p{0.14\linewidth} p{0.14\linewidth} p{0.14\linewidth} p{0.3\linewidth}}
\toprule
\textbf{Parameter} & \textbf{Stage 0-1 (Stable)} & \textbf{Stage 2-3 (Mild Chaos)} & \textbf{Stage 4-6 (Chaotic)} & \textbf{Rationale} \\
\midrule
SIREN first\_omega\_0 & 30.0 & 20.0 & 10.0 & Matches network frequency resolution to data complexity. \\
Energy Weight $\lambda_E$ & $10^{-3}$ - $10^{-2}$ & $10^{-1}$ - $1$ & $10^{1}$ - $10^{3}$ & Increases energy conservation constraint as chaos increases. \\
Adam Initial LR & $10^{-3}$ - $7\text{e-}4$ & $4.9\text{e-}4$ - $3.4\text{e-}4$ & $2.4\text{e-}4$ - $1.2\text{e-}4$ & Gradually reduces initial learning rate for more difficult stages. \\
Gradient Clip $C$ & 8.5 & 5.0 & 3.0 & Allows larger gradients for stable data; clips aggressively for chaotic data. \\
L-BFGS Type & Full-Batch & Full-Batch / sL-BFGS & sL-BFGS & Switches to stochastic variant as cumulative dataset size grows. \\
ALM State & Reset & Reset & Reset & Resets ALM statistics to adapt to new data distribution at each stage. \\
\bottomrule
\end{tabular}
\renewcommand{\arraystretch}{1.0}
\end{table}
